\setcounter{framenumber}{0}
%24

\begin{frame}
  \titlepage
\end{frame}

\begin{frame}{Learning Goals}
By the end of this lecture, students should be able to:
\begin{itemize}
    \item define the Chi-Square distribution;
    \item understand the connection between Chi-Square and Gamma distributions;
    \item compute means and variances of Chi-Square random variables;
    \item understand why Chi-Square distributions appear in statistics;
    \item define the Student \(t\)-distribution;
    \item understand why \(t\)-distributions have heavier tails than Normal distributions;
    \item understand convergence of the \(t\)-distribution to the Normal distribution.
\end{itemize}
\end{frame}

\lecturepart{Chi-Square Distribution}

\begin{frame}{Why Chi-Square Distributions Matter}
Chi-Square distributions appear throughout statistics:
\begin{itemize}
    \item sample variance;
    \item hypothesis testing;
    \item confidence intervals;
    \item goodness-of-fit tests;
    \item contingency tables.
\end{itemize}

\vspace{0.25cm}

\begin{keybox}{Main idea}
Chi-Square distributions measure accumulated squared Normal variation.
\end{keybox}
\end{frame}

\begin{frame}{Definition of the Chi-Square Distribution}
\begin{keybox}{Definition}
Suppose
\[
Z_1,\dots,Z_n
\]
are independent standard Normal random variables.

Define
\[
X
=
Z_1^2+\cdots+Z_n^2.
\]

Then
\[
\boxed{
X\sim \chi_n^2.
}
\]

We say that \(X\) has a Chi-Square distribution with
\[
n
\]
degrees of freedom.
\end{keybox}
\end{frame}

\begin{frame}{Why Squaring Matters}
Each
\[
Z_i^2
\]
measures squared deviation from the mean.

\vspace{0.25cm}

Adding them gives a total amount of variation:
\[
Z_1^2+\cdots+Z_n^2.
\]

\vspace{0.25cm}

\begin{keybox}{Interpretation}
Chi-Square random variables measure total squared fluctuation.
\end{keybox}
\end{frame}

\begin{frame}{Chi-Square and Gamma}
Recall:
if
\[
X\sim \operatorname{Gamma}(\alpha,\lambda),
\]
then
\[
f(x)
=
\frac{
\lambda^\alpha
}{
\Gamma(\alpha)
}
x^{\alpha-1}e^{-\lambda x}.
\]

\vspace{0.25cm}

A Chi-Square distribution is a special Gamma distribution:
\[
\boxed{
\chi_n^2
=
\operatorname{Gamma}
\left(
\frac n2,
\frac12
\right).
}
\]
\end{frame}

\begin{frame}{PDF of the Chi-Square Distribution}
Using the Gamma formula:
\[
\boxed{
f(x)
=
\frac{
1
}{
2^{n/2}\Gamma(n/2)
}
x^{n/2-1}e^{-x/2},
\qquad x>0.
}
\]

\vspace{0.25cm}

The shape depends heavily on the degrees of freedom.
\end{frame}

\begin{frame}{Mean and Variance of Chi-Square}
If
\[
X\sim \chi_n^2,
\]
then
\[
\boxed{
\mathbb{E}(X)=n.
}
\]

Also,
\[
\boxed{
\operatorname{Var}(X)=2n.
}
\]

\vspace{0.25cm}

\begin{keybox}{Observation}
Both the mean and variance grow linearly with the degrees of freedom.
\end{keybox}
\end{frame}

\begin{frame}{Shape of Chi-Square Distributions}
Small degrees of freedom:
\begin{itemize}
    \item strongly right-skewed;
    \item concentrated near zero.
\end{itemize}

\vspace{0.25cm}

Large degrees of freedom:
\begin{itemize}
    \item less skewed;
    \item more symmetric;
    \item approximately Normal.
\end{itemize}

\vspace{0.25cm}

\begin{keybox}{Connection to CLT}
A Chi-Square distribution becomes approximately Normal for large
\[
n.
\]
\end{keybox}
\end{frame}

\begin{frame}{Sample Variance and Chi-Square}
Suppose
\[
X_1,\dots,X_n
\overset{\text{i.i.d.}}{\sim}
N(\mu,\sigma^2).
\]

Recall the sample variance:
\[
S_n^2
=
\frac1{n-1}
\sum_{j=1}^n
(X_j-\bar X_n)^2.
\]

Then:
\[
\boxed{
\frac{
(n-1)S_n^2
}{
\sigma^2
}
\sim
\chi_{n-1}^2.
}
\]

\vspace{0.25cm}

This is one of the most important results in classical statistics.
\end{frame}

\begin{frame}{Degrees of Freedom}
Why do we get
\[
n-1
\]
degrees of freedom instead of
\[
n?
\]

\vspace{0.25cm}

Because:
\[
\bar X_n
\]
has already been estimated from the data.

\vspace{0.25cm}

Once the sample mean is fixed,
only
\[
n-1
\]
deviations can vary independently.
\end{frame}

\begin{frame}{Chi-Square Additivity}
\begin{keybox}{Theorem}
If
\[
X\sim \chi_m^2,
\qquad
Y\sim \chi_n^2,
\]
independently, then
\[
\boxed{
X+Y
\sim
\chi_{m+n}^2.
}
\]
\end{keybox}

\vspace{0.25cm}

This follows immediately from the Gamma representation.
\end{frame}

\lecturepart{Student \(t\)-Distribution}

\begin{frame}{Why the \(t\)-Distribution?}
Suppose:
\begin{itemize}
    \item the population is Normal;
    \item the variance
    \[
    \sigma^2
    \]
    is unknown;
    \item the sample size is small.
\end{itemize}

\vspace{0.25cm}

Then the standard Normal approximation is no longer exact.

\vspace{0.25cm}

\begin{keybox}{Solution}
The Student \(t\)-distribution accounts for uncertainty in estimating the variance.
\end{keybox}
\end{frame}

\begin{frame}{Definition of the Student \(t\)-Distribution}
\begin{keybox}{Definition}
Let
\[
Z\sim N(0,1),
\]
and
\[
V\sim \chi_n^2,
\]
independently.

Define
\[
\boxed{
T
=
\frac{
Z
}{
\sqrt{V/n}
}.
}
\]

Then
\[
T
\]
has the Student \(t\)-distribution with
\[
n
\]
degrees of freedom.
\end{keybox}
\end{frame}

\begin{frame}{Why the \(t\)-Distribution Has Heavy Tails}
The denominator:
\[
\sqrt{V/n}
\]
is random.

\vspace{0.25cm}

Sometimes it becomes unusually small.

\vspace{0.25cm}

That makes:
\[
|T|
\]
occasionally very large.

\vspace{0.25cm}

\begin{keybox}{Consequence}
The \(t\)-distribution has heavier tails than the Normal distribution.
\end{keybox}
\end{frame}

\begin{frame}{Properties of the \(t\)-Distribution}
The \(t\)-distribution is:
\begin{itemize}
    \item symmetric about
    \[
    0;
    \]

    \item bell-shaped;

    \item heavier-tailed than the Normal distribution.
\end{itemize}

\vspace{0.25cm}

As the degrees of freedom increase:
\begin{itemize}
    \item tails become lighter;
    \item the distribution approaches the Normal distribution.
\end{itemize}
\end{frame}

\begin{frame}{Convergence to the Normal Distribution}
As
\[
n\to\infty,
\]
the Chi-Square random variable satisfies:
\[
\frac{V}{n}\to1.
\]

Therefore:
\[
T
=
\frac{
Z
}{
\sqrt{V/n}
}
\approx Z.
\]

Thus:
\[
\boxed{
t_n
\to
N(0,1).
}
\]
\end{frame}

\begin{frame}{The Cauchy Distribution}
If
\[
n=1,
\]
then the \(t\)-distribution becomes the Cauchy distribution.

\vspace{0.25cm}

The Cauchy distribution has:
\begin{itemize}
    \item extremely heavy tails;
    \item undefined mean;
    \item undefined variance.
\end{itemize}

\vspace{0.25cm}

\begin{keybox}{Important}
Not every distribution has finite expectation or variance.
\end{keybox}
\end{frame}

\begin{frame}{Confidence Intervals}
The \(t\)-distribution is fundamental for confidence intervals.

\vspace{0.25cm}

For example:
\[
\bar X_n
\pm
t^\ast
\frac{
S_n
}{
\sqrt n
}.
\]

\vspace{0.25cm}

Here:
\begin{itemize}
    \item
    \[
    \bar X_n
    \]
    is the sample mean;

    \item
    \[
    S_n
    \]
    estimates the standard deviation;

    \item
    \[
    t^\ast
    \]
    comes from the \(t\)-distribution.
\end{itemize}
\end{frame}

\begin{frame}{Normal vs \(t\)-Distribution}
\begin{center}
\renewcommand{\arraystretch}{1.5}
\begin{tabular}{lll}
\toprule
Distribution & Tails & Variance known? \\
\midrule
Normal & lighter & yes \\
Student \(t\) & heavier & no \\
\bottomrule
\end{tabular}
\end{center}

\vspace{0.25cm}

\begin{keybox}{Interpretation}
The heavier tails account for uncertainty in estimating the variance.
\end{keybox}
\end{frame}

\begin{frame}{Big Picture}
Three major distributions now fit together:

\vspace{0.25cm}

\begin{center}
\renewcommand{\arraystretch}{1.6}
\begin{tabular}{ccc}
\toprule
Distribution & Built From & Main Role \\
\midrule
Normal & sums/CLT & averages \\
Chi-Square & squared Normals & variances \\
Student \(t\) & Normal divided by Chi-Square & inference \\
\bottomrule
\end{tabular}
\end{center}
\end{frame}

\begin{frame}{Summary}
\begin{itemize}
    \item Chi-Square distributions are sums of squared standard Normals.

    \item Chi-Square distributions are special Gamma distributions.

    \item Sample variances naturally produce Chi-Square distributions.

    \item Student \(t\)-distributions arise when variance is estimated.

    \item \(t\)-distributions have heavier tails than Normals.

    \item As degrees of freedom increase:
    \[
    t_n
    \to
    N(0,1).
    \]
\end{itemize}

\vspace{0.25cm}

\begin{keybox}{End of Probability I}
The LLN, CLT, Chi-Square, and \(t\)-distributions form the mathematical foundation of classical statistics.
\end{keybox}
\end{frame}